\documentclass[11pt,a4paper]{article}
\usepackage[margin=24mm]{geometry}
\usepackage{amsmath,amssymb,amsthm,booktabs,array,enumitem,hyperref,fancyhdr}
\hypersetup{colorlinks=true,linkcolor=blue,urlcolor=blue,citecolor=blue}
\newtheorem{theorem}{Theorem}
\newtheorem{proposition}[theorem]{Proposition}
\newtheorem{corollary}[theorem]{Corollary}
\newcommand{\C}{\mathbb C}
\newcommand{\R}{\mathbb R}
\newcommand{\U}{\mathrm U}
\newcommand{\tr}{\operatorname{tr}}
\newcommand{\rank}{\operatorname{rank}}
\newcommand{\Id}{I}
\newcommand{\ad}{\dagger}
\pagestyle{fancy}\fancyhf{}\fancyhead[L]{UD gauge transport / development v0.1}
\fancyhead[R]{15 October 2026 UTC}\fancyfoot[C]{\thepage}
\setlength{\headheight}{14pt}
\setlength{\parskip}{5pt}\setlength{\parindent}{0pt}
\setlist{itemsep=3pt,topsep=4pt}
\title{Gauge-Covariant Incidence Transport on a Filled Tetrahedron\\
\large A flatness theorem and an exact phase-controlled spectrum}
\author{Benjamin Walker Mayes\\\small AI-assisted development manuscript}
\date{15 October 2026 UTC / version 0.1}
\begin{document}
\maketitle
\begin{center}\small
\textbf{EXACT CONDITIONAL mathematics. CANDIDATE gauge extension.}\\
Constructor verification passed; independent review pending. Physical promotion: \textbf{0}.
\end{center}
\begin{abstract}
We construct a finite unitary transport extension of the declared tetrahedral
incidence flow. A rank-$r$ complex fiber is assigned to each simplex and an
independent unitary map to each adjacent incidence. The resulting generator is
skew-Hermitian, has gauge-invariant support currents, and is covariant under
local frame changes and signed tetrahedral relabelings. Its degree-two square
is the signed difference of transport along the two sides of an incidence
diamond. On the filled tetrahedron, preserving the chain identity forces all
such curvature to vanish and the connection to be pure gauge. Curved
transport remains norm conserving but changes the original Hodge-square
identity. A single phase on one vertex--edge incidence gives an exactly
solvable spectral splitting, providing a finite analogue test contract.
The construction supplies a precise candidate interface, not a derivation
of the Standard Model, a gauge-field action, a continuum limit, running
couplings or physical unification. The underlying connection/curvature
principles are established mathematics; no general novelty claim is made.
\end{abstract}

\section{Source commitments and added objects}
The source carrier is the downward-closed filled tetrahedron $\Delta^3$,
with increasingly oriented simplices and graded dimensions $(4,6,4,1)$.
Let $B$ be the total boundary, acting downward on the direct sum of all
15 coordinates. The native real flow is
\[
 B^2=0,\qquad A_0=B^T-B,\qquad \dot c=A_0c .
\]
Its complexification and unit simplex metric are mathematical choices,
not a physical quantization. Papers I--III, review drafts 1.1.1, retain
their source identities \cite{UD1,UD2,UD3}. In particular, the finite
$S_4$ symmetry is not itself a continuous physical gauge group.

The Hasse graph has one node per nonempty simplex and one link for each
adjacent containment. It has 15 nodes and 28 links. We add the following:
\begin{center}\small
\begin{tabular}{p{.24\linewidth}p{.67\linewidth}}
\toprule
Object & Type and status\\ \midrule
$E_\sigma=\C^r$ & Added internal fiber at simplex $\sigma$; $r$ is supplied.\\
$c_\sigma\in E_\sigma$ & Complex amplitude vector, with support $q_\sigma=\|c_\sigma\|^2$.\\
$U_{\sigma\tau}$ & Unitary map $E_\tau\to E_\sigma$ for $\sigma\subset\tau$ of codimension one. Reverse transport is its adjoint.\\
$s$ & Internal evolution parameter; no laboratory clock is assigned.\\
$\Omega_\sigma(s)$ & Optional local skew-Hermitian temporal connection.\\
\bottomrule
\end{tabular}
\end{center}
The fibers are indexed by simplices, not merely by the four geometric
vertices. Thus the local frame group is $\prod_\sigma\U(r)$.
This is an explicit extension with additional state data. It is not a
discovery that the original scalar coordinate already carries color,
weak isospin or hypercharge.

Define a downward block operator by
\[
 (D_U)_{\sigma\tau}=B_{\sigma\tau}U_{\sigma\tau},\qquad
 A_U=D_U^\ad-D_U,\qquad
 L_U=D_U^\ad D_U+D_UD_U^\ad.                         \tag{1}
\]
Unlisted blocks vanish. For identity transport, $D_U=B\otimes I_r$.
The original topological boundary $B$ remains defined even when the
new transport operator $D_U$ is not nilpotent.

\section{The transport and flatness theorem}
A codimension-two containment $\sigma\subset\eta$ has exactly two
intermediate simplices $\tau,\tau'$. Write
$\epsilon=B_{\sigma\tau}B_{\tau\eta}$; the other path has sign $-\epsilon$.
Its transport defect and based diamond holonomy are
\[
 F_{\sigma\eta}
 =U_{\sigma\tau}U_{\tau\eta}
  -U_{\sigma\tau'}U_{\tau'\eta},\qquad
 W=(U_{\sigma\tau'}U_{\tau'\eta})^\ad
       U_{\sigma\tau}U_{\tau\eta}.                   \tag{2}
\]
Both are independent of the matter state. $F$ maps $E_\eta$ to $E_\sigma$;
$W$ is an endomorphism of $E_\eta$.

\begin{theorem}[Finite gauge-transport admissibility]\label{th:main}
On the filled tetrahedron, for every integer $r\geq1$ and every assignment
of unitary incidence transports, the following statements hold.
\begin{enumerate}
\item $A_U^\ad=-A_U$. For continuous prescribed links, the equation
$\dot c=(A_U+\Omega)c$ preserves total support when
$\Omega=\bigoplus_\sigma\Omega_\sigma$ is skew-Hermitian.
\item Under $c'_\sigma=g_\sigma c_\sigma$ and
$U'_{\sigma\tau}=g_\sigma U_{\sigma\tau}g_\tau^\ad$,
$D_{U'}=GD_UG^\ad$ and $A_{U'}=GA_UG^\ad$.
For differentiable time-dependent frames, full evolution covariance
requires $\Omega'=G\Omega G^\ad+\dot G G^\ad$.
\item Every codimension-two block is
$(D_U^2)_{\sigma\eta}=\epsilon F_{\sigma\eta}$.
Consequently $D_U^2=0$ if and only if every diamond holonomy is identity.
\item If $D_U^2=0$, there are unitary $P_\sigma$ such that
$U_{\sigma\tau}=P_\sigma P_\tau^\ad$. Therefore $A_U$ is unitarily
conjugate to $A_0\otimes I_r$ and
\[
 \det(\lambda I-A_U)=\lambda^r(\lambda^2+4)^{7r}.     \tag{3}
\]
\item For arbitrary links,
\[
 -A_U^2=L_U-D_U^2-(D_U^\ad)^2.                     \tag{4}
\]
Thus the unchanged degree-diagonal Hodge-square identity $-A_U^2=L_U$
holds if and only if the connection is flat in the sense of item 3.
\end{enumerate}
\end{theorem}

\begin{proof}
Adjoint transposition of (1) proves skew-Hermiticity. With
$Q=\sum_\sigma c_\sigma^\ad c_\sigma$, differentiation gives
$\dot Q=c^\ad[(A_U+\Omega)^\ad+(A_U+\Omega)]c=0$.
Each nonzero block transforms by the stated left/right frame factors;
summing proves covariance. Differentiating $c'=Gc$ supplies the
inhomogeneous $\dot GG^\ad$ term. Omitting it is valid only for
time-independent frames, or a separately fixed temporal gauge.

For item 3, a two-step downward path from $\eta$ to $\sigma$ removes
two distinct vertices. Exactly two orders are possible, and the
ordinary incidence signs are opposite. Their two transport products
are therefore precisely $\epsilon F_{\sigma\eta}$. All other blocks
of $D_U^2$ vanish by degree. Since the second path product is unitary,
$F=0$ is equivalent to $W=I$.

For item 4 let $t$ denote the unique top tetrahedron. Choose
$P_t=I$ and define $P_\sigma:E_t\to E_\sigma$ by transport along a
downward path from $t$ to $\sigma$, without its incidence signs.
Such paths correspond to orders of deleting the vertices absent
from $\sigma$. Adjacent transpositions generate all changes of order;
each transposition is exactly a diamond. Flatness makes the path
product independent of the order. In particular
$P_\sigma=U_{\sigma\tau}P_\tau$ for each incidence. All $P_\sigma$
are unitary, so
\[
 D_U=P(B\otimes I_r)P^\ad,\qquad
 P=\bigoplus_\sigma P_\sigma .
\]
This proves the conjugacy and (3), using the native spectrum.
The unique top cell is substantive: flat connections on a
non-simply-connected base need not be globally pure gauge.

Finally expand $(D_U^\ad-D_U)^2$ to obtain (4). The two correction
terms lower and raise degree by two respectively. They occupy
disjoint graded blocks, so their sum vanishes only when each vanishes.
This proves item 5.
\end{proof}

The flatness conclusion concerns incidence transport at each fixed $s$.
It does not exclude mixed temporal curvature
\[
 \mathcal E_{\sigma\tau}
 =\dot U_{\sigma\tau}-\Omega_\sigma U_{\sigma\tau}
                    +U_{\sigma\tau}\Omega_\tau ,
\]
which transforms as $g_\sigma\mathcal E_{\sigma\tau}g_\tau^\ad$.
For example, identity links with unequal local scalar temporal blocks
can have $\mathcal E\ne0$ while $D_U^2=0$.
Instantaneous conjugacy of $A_U$ therefore does not establish equality
of the complete time-dependent propagators. The spectral corollary
below fixes $\Omega=0$ and constant links.

\subsection{Support currents and full tetrahedral covariance}
For each adjacent incidence put
\[
 J_{\sigma\tau}
 =2\operatorname{Re}
     (c_\sigma^\ad B_{\sigma\tau}U_{\sigma\tau}c_\tau).       \tag{5}
\]
Its contribution is $-J_{\sigma\tau}$ to $\dot q_\sigma$ and
$+J_{\sigma\tau}$ to $\dot q_\tau$. Local temporal blocks contribute
zero to each $q_\sigma$. Thus for any coordinate region $R$,
\[
 \frac{d}{ds}\sum_{\sigma\in R}q_\sigma
 =J_{\rm in}(R)-J_{\rm out}(R).
\]
Equation (5) is gauge invariant because both endpoint frame factors
cancel. It is a model support current, not an electrical-current
identification.

For $\pi\in S_4$ let $S_\pi$ be the signed permutation of oriented
simplices. Move the independent link data by
$U^\pi_{\pi\sigma,\pi\tau}=U_{\sigma\tau}$.
Incidence equivariance gives
\[
 D_{U^\pi}=(S_\pi\otimes I_r)D_U(S_\pi^T\otimes I_r).       \tag{6}
\]
This is covariance of the whole connection family. A particular
nonuniform connection need not commute with $S_\pi$. No curved
configuration is declared tetrahedrally invariant merely because
the law is covariant.

For a diamond a dimensionless gauge-invariant curvature diagnostic is
\[
 \chi_W=1-\frac1r\operatorname{Re}\tr W
       =\frac{\|F_{\sigma\eta}\|_F^2}{2r}\in[0,2].           \tag{7}
\]
This follows by expanding the Frobenius norm of the difference of
the two unitary path products. On the scalar graph, a spanning tree
removes 14 of the 28 link phases, leaving 14 independent loop
coordinates. The 18 diamond equations generate all loop constraints
here by the path proof above. The replay also checks their
linearized scalar rank, 14. These are model coordinate counts,
not physical particle or spacetime dimensions.

\section{An exact, phase-controlled spectral prediction}
Set $r=1$, $\Omega=0$, all links to one except
$U_{\{0\},\{0,1\}}=u=e^{i\phi}$. Keep every link magnitude and every
incidence sign fixed.

\begin{corollary}[One-incidence phase spectrum]\label{cor:phase}
For this declared family,
\[
 \det(\lambda I-A_U)
 =\lambda(\lambda^2+4)^5
   \left[(\lambda^2+4)^2-4(1-\cos\phi)\right].              \tag{8}
\]
The positive squared frequencies are $4$ five times and
$4\pm2\sqrt{1-\cos\phi}$ once each, with multiplicities combined
at $\phi=0$. At $u=i$,
\[
 \det(\lambda I-A_U)
 =\lambda(\lambda^2+2)(\lambda^2+4)^5(\lambda^2+6).         \tag{9}
\]
The connection is curved for $u\ne1$ even though all support balances
remain exact. At $u=i$, $\rank D_U^2=1$ and
$\rank(-A_U^2-L_U)=2$.
\end{corollary}

\begin{proof}
Partition the 15 coordinates into eight even-degree and seven
odd-degree coordinates. In this order write
$A_U=\left(\begin{smallmatrix}0&C\\-C^\ad&0\end{smallmatrix}\right)$.
For the identity-link matrix $C_0$, the native identity gives
$C_0^\ad C_0=4I_7$. The vertex-0 row has three unit entries.
Let $e$ select odd coordinate edge 01 and let
$v=C_0^\ad e_{\text{vertex }0}=e+w$, so $w^\ad e=0$ and
$\|w\|^2=2$. With $\delta=u-1$,
\[
 C=C_0+\delta e_{\text{vertex }0}e^\ad,\qquad
 C^\ad C=4I_7+\delta we^\ad+\bar\delta ew^\ad .
\]
The diagonal $ee^\ad$ correction cancels since
$\delta+\bar\delta+|\delta|^2=0$. On the plane spanned by
$e,w/\sqrt2$, the added Hermitian matrix has eigenvalues
$\pm\sqrt2|\delta|=\pm2\sqrt{1-\cos\phi}$.
It is zero on the other five directions. The nonzero eigenvalues
of the bipartite skew matrix are the signed imaginary singular
values of $C$, with one additional zero direction. This proves (8).
For $u\ne1$, $D_U^2$ has only the row at vertex 0, with nonzero
entries from faces 012 and 013, so its rank is one. Its adjoint
occupies the opposite graded block, proving the defect rank.
\end{proof}

The corollary is a finite-model testable prediction, not an
electroweak or strong-coupling spectrum. Frequencies and their ratios
are insensitive to a common positive clock multiplier, but an
actual sensor must still be calibrated.

\section{What the native symmetries do and do not select}
\begin{proposition}[Global symmetry versus local gauge data]
The real orthogonal centralizer of the canonical native $A_0$ is
$O(1)\times\U(7)$. Its subgroup also commuting with the full signed
$S_4$ action is $O(1)\times\U(1)^3$. Consequently the connected
$S_4$-compatible global symmetry has no nontrivial continuous
homomorphism from $SU(2)$ or $SU(3)$.
\end{proposition}
\begin{proof}
The harmonic line is $\ker A_0$ and must be preserved. On its
orthogonal complement $J=A_0/2$ is an orthogonal complex structure.
Orthogonal maps commuting with $J$ are exactly $\U(7)$.
As a complex $S_4$ representation this space is the sum of one
copy each of the inequivalent real-type irreducibles of complex
dimensions $1,3,3$. Schur's lemma gives one unit phase per summand.
A homomorphism from either connected simple group to the abelian
torus has zero Lie algebra map, hence trivial image.
\end{proof}
This result does not exclude differently embedded, noncommuting
or enlarged gauge branches. In particular $\U(7)$ contains
non-Abelian subgroups if the simultaneous $S_4$-commuting restriction
is dropped. It supplies neither a chosen embedding nor local
parallel transport. The requirement of commuting with $S_4$ is
an explicit branch assumption, not a universal law.

For a separate comparison, exact scalar state phases
$z_\sigma=c_\sigma/|c_\sigma|$ on nonzero support produce receipts
$u_{\sigma\tau}=\bar z_\sigma z_\tau$. Their product around every
closed support-graph path is one by telescoping, with conjugation
on reversed edges. Zeros remove links and require separate
components. Therefore these relative-state receipts cannot alone
instantiate an independent nonflat connection. They remain valid
reconstruction data as in Paper II; no source result is withdrawn.

\section{Coupling selection remains a separate theorem target}
Even after \emph{assuming} the Standard Model Lie algebra
$\mathfrak{su}(3)\oplus\mathfrak{su}(2)\oplus\mathfrak u(1)$, gauge
invariance does not impose a common kinetic coefficient.
Every positive invariant symmetric bilinear form is
\[
 h=a_3h_3\oplus a_2h_2\oplus a_1h_1,\qquad a_i>0,          \tag{10}
\]
once generator normalizations are fixed.
On each compact simple ideal the adjoint representation and
Schur's lemma give one scalar multiple of a fixed invariant form.
Cross terms vanish because each simple ideal is perfect and
commutes with the other ideals and the center. The one-dimensional
center has its independent positive coefficient. This proves
the three-parameter freedom without fitting a numerical target.

A chosen compact simple algebra instead has one such form up
to scale. Conditional $SU(5)$ illustrates the normalization
obligation: if one supplies
\[
 Y=\operatorname{diag}(-1/3,-1/3,-1/3,1/2,1/2),\quad
 \tr Y^2=5/6,\quad T_Y=\sqrt{3/5}\,Y ,
\]
then $\tr T_Y^2=1/2$ and a common $g_5T_Y=g_YY$ gives
$g_1=\sqrt{5/3}\,g_Y=g_5$ at tree-level matching.
This is standard conditional group theory \cite{PDG}, not selection
of $SU(5)$ by UD and not a running-coupling calculation. A
renormalization scheme, matter spectrum, thresholds, action and
quantization remain required for scale evolution.

\section{A positive candidate contract and its exit gates}
The candidate branch has complete finite \emph{matter} kinematics:
choose a fixed tetrahedron, rank $r$, unitary links $U(s)$, local
temporal blocks $\Omega(s)$, initial $c(0)$ and readouts (5),(7).
It has well-defined norm-preserving linear evolution and full
family relabeling covariance. On the curved branch $D_U$ is a
covariant transport operator, not the original topological boundary.

The independent link variables currently have no source-selected
evolution. A nonnegative gauge-invariant potential can, for example,
be \emph{proposed} as
$V=\sum_d\kappa_d[1-\operatorname{Re}\tr W_d/r]$, $\kappa_d>0$.
This is a supplied Wilson-type choice, not a derived gauge action.
A potential alone does not define a link clock, canonical momenta,
Gauss constraint, feedback or quantum measure. Symmetry may constrain
some coefficient classes without fixing their values.

\begin{center}\small
\begin{tabular}{p{.35\linewidth}p{.56\linewidth}}
\toprule
Admissibility requirement & Outcome for this branch\\ \midrule
Explicit states, fibers and transport & Defined, conditional on supplied $r,U,\Omega$.\\
Local gauge and signed $S_4$ family covariance & Proved; configuration invariance is distinct.\\
Conserved support and regional current law & Proved, including curved transport.\\
Original chain identity and Hodge square & Preserved only on the flat branch; proved scope.\\
Independent curvature readout & Defined; acquisition and physical interpretation open.\\
Closed gauge-field dynamics and Gauss law & Open; needs a joint action or phase-space law.\\
Continuum/Lorentz and chiral matter limit & Open; no convergence transferred from algebra.\\
SM group, charges, anomalies, normalization & Open; no native selection is supplied.\\
RG running, proton channels, monopole sector & Downstream of the missing field/action bridge.\\
\bottomrule
\end{tabular}
\end{center}
Discrete connections and the curvature/nilpotency relationship are
established \cite{BEHS}. Algebraic covariance alone is not a convergence
theorem: other discretizations specifically investigate convergence
\cite{Braune}. Our Hasse-link construction is not identified with
their smooth discretizations or their continuum limit.

\subsection{Conditional analogue experiment specification}
A bounded operator analogue of Corollary \ref{cor:phase} would require:
15 calibrated complex modes, the 28 prescribed signed incidence
couplings, matched coupling magnitudes, a controlled phase on the
chosen incidence and a known common rate $\eta>0$. In a proposed
Schrodinger-type realization use $H_{\rm phys}=i\hbar\eta A_U$;
then the evolution is $\dot c=\eta A_Uc$ in laboratory time.
At phases $0$ and $\pi/2$, the predicted nonzero frequency sets are
respectively $\{2\eta\}$ and
$\{\sqrt2\eta,2\eta,\sqrt6\eta\}$ with the multiplicities in (9).
Measured complex transfer matrices or several calibrated preparations
are preferable to a single power-spectrum trace that may miss modes.

Before testing, freeze the mapping, phase uncertainty, hopping errors,
readout noise/covariance, frequency estimator, time window, null model,
acceptance threshold, exclusions and retained data. Gauge-related
rephasings are the equivalence control; uniform magnitude changes
are a distinct confound. This is a model-to-engineered-operator
candidate, not a test of GUT unification or proton/monopole physics.
No physical experiment or dataset fit has been performed.

\section{Verification, provenance and review}
The fresh primary replay \texttt{verify\_hasse\_transport.py} constructs
all simplices and incidence signs, enumerates 18 diamonds, checks
the 14-dimensional scalar cycle constraint, tests non-Abelian
rank-two covariance, all 24 signed relabelings, temporal covariance,
local support balance and the phase-dependent characteristic polynomial
on five exact phases. The spectral proof above uses a separate
two-plane Gram reduction, not the computed polynomial as a universal
proof. The companion \texttt{verify\_gauge\_bridge.py} independently
checks native centralizer dimensions, state-phase flatness, conditional
$SU(5)$ trace normalization and an ordered-vertex transport comparison.
That comparison is an additional convention, not the Hasse-link
model and not a claim of its full permutation covariance.

The receipts record exact rational, Gaussian-rational and algebraic
arithmetic. They are constructor checks, not independent review.
No source script was executed merely to ingest an archive.
Source baseline: Public Atlas commit
\texttt{84d6636cbf7d5dcb3c8c9e8abbe7b43350908499};
UD repository baseline
\texttt{2a6d3fda46dd4724aeba73329b605e6533d68d5d}.
The historical v0.39 ledger is preserved alongside later source
amendments. The new objects and conditional conclusions do not
change frozen physical status or claim confidence by citation.
The unavailable Genesis image is not used.

Author responsibility and independent mathematical review remain
submission tasks. AI assistance includes proof drafting, exact
fixtures, source comparison and document preparation. No claim
of new general gauge theory, funding status, experimental discovery
or physical promotion is made.

\begin{thebibliography}{9}
\bibitem{UD1} B. W. Mayes, \emph{Finite Tetrahedral Complexes, Skew Support
Flow and Contact Readiness}, Paper I, review draft 1.1.1, 4 October 2026,
Sections 4--6, Public Research Atlas.
\bibitem{UD2} B. W. Mayes, \emph{Projective Reconstruction, Complex Phase
Receipts and Frozen-Tree Validation}, Paper II, review draft 1.1.1,
4 October 2026, Sections 4 and 6.
\bibitem{UD3} B. W. Mayes, \emph{Tetrahedral Symmetry, Radial Spectra and
Finite Collective Mechanisms}, Paper III, review draft 1.1.1,
4 October 2026, Section 2.3.
\bibitem{BEHS} D. Berwick-Evans, A. N. Hirani and M. D. Schubel,
\emph{Discrete Vector Bundles with Connection},
\href{https://arxiv.org/html/2104.10277v3}{arXiv:2104.10277v3},
Sections 4--6; inspected 15 October 2026 UTC.
\bibitem{Braune} T. Braune, Y. Tong, F. Gay-Balmaz and M. Desbrun,
\emph{A Discrete Exterior Calculus of Bundle-valued Forms},
\href{https://arxiv.org/html/2406.05383v1}{arXiv:2406.05383v1},
Sections 4--5; inspected 15 October 2026 UTC.
\bibitem{PDG} N. Nagata, A. Hebecker and J. Hisano,
\emph{Grand Unified Theories}, 2025 Review of Particle Physics,
Sections 92.2.2 and 92.5,
\href{https://pdg.lbl.gov/2025/reviews/rpp2025-rev-guts.pdf}{PDG review}.
\end{thebibliography}
\end{document}
